\section{Conclusion}
The seminar paper is a discussion of the challenges of resource sharing in real-time systems, particularly focusing on priority inversion and deadlocks. It has been discussed that how these issues can lead to significant operational problems, including missed deadlines and system failures. The chapters also provided idea of classic protocols designed to manage resource sharing effectively, such as Non-Preemptive Protocol (NPP), Priority Inheritance Protocol (PIP), Priority Ceiling Protocol (PCP), and Stack Resource Policy (SRP). Each protocol has its own benefits and drawbacks, making them suitable for different scenarios in real-time systems based upon operations.
The discussions highlighted the importance of selecting the appropriate protocol based on the specific requirements and constraints of the given system. By understanding these protocols, developers can better design real-time systems that are robust by ensuring timely and reliable operations.

